% !TEX program = lualatex
\documentclass[final]{beamer}

% Gemini-style portrait poster starter
% 24in x 36in portrait, 3 columns
% Compile with: pdflatex (or lualatex/xelatex if you add custom fonts)

\usepackage[size=custom,width=60.96,height=91.44,scale=1.0]{beamerposter}
\usetheme{gemini}

\usepackage{graphicx}
\usepackage{booktabs}
\usepackage{amsmath, amssymb, bm}
\usepackage{multicol}
\usepackage{array}
\usepackage{ragged2e}
\usepackage{caption}
\usepackage{xcolor}
\usepackage{tikz}
% enumitem is incompatible with beamer; use beamer's native list spacing instead

% -----------------------------
% Gemini-like color palette
% -----------------------------
\definecolor{GeminiBlue}{RGB}{33, 97, 140}
\definecolor{GeminiLightBlue}{RGB}{236, 244, 252}
\definecolor{GeminiDark}{RGB}{36, 41, 46}
\definecolor{GeminiGray}{RGB}{90, 98, 108}
\definecolor{GeminiAccent}{RGB}{0, 114, 178}

\setbeamercolor{headline}{fg=white,bg=GeminiBlue}
\setbeamercolor{title in headline}{fg=white}
\setbeamercolor{author in headline}{fg=white}
\setbeamercolor{institute in headline}{fg=white}

\setbeamercolor{block title}{fg=GeminiBlue,bg=white}
\setbeamercolor{block body}{fg=black,bg=white}

\setbeamercolor{block alerted title}{fg=white,bg=GeminiBlue}
\setbeamercolor{block alerted body}{fg=black,bg=GeminiLightBlue}

\setbeamercolor{block example title}{fg=white,bg=GeminiAccent}
\setbeamercolor{block example body}{fg=black,bg=GeminiLightBlue}

\setbeamercolor{normal text}{fg=GeminiDark,bg=white}
\setbeamercolor{footline}{fg=white,bg=GeminiBlue}

% -----------------------------
% Layout
% -----------------------------
\newlength{\sepwidth}
\newlength{\colwidth}
\setlength{\sepwidth}{0.025\paperwidth}
\setlength{\colwidth}{0.30\paperwidth}
% 1 sep + 3 cols + 2 inner seps + 1 sep = 2*sep + 3*col + 2*sep = 4*sep + 3*col
% Here: 4*0.025 + 3*0.30 = 1.00

\setlength{\parskip}{0.4em}
\renewcommand{\arraystretch}{1.2}

\captionsetup{font=large}
\setlength{\leftmargini}{1.5em}
\setlength{\leftmarginii}{1.5em}

% -----------------------------
% Metadata
% -----------------------------
\title{Your Poster Title Goes Here}
\author{First Author \and Second Author \and Third Author}
\institute{Department or Lab \quad\textbar\quad University or Organization}

% Optional logos
% \logoright{\includegraphics[height=5cm]{logo_right.pdf}}
% \logoleft{\includegraphics[height=5cm]{logo_left.pdf}}

% -----------------------------
% Helper commands
% -----------------------------
\newcommand{\panelspace}{\vspace{0.6em}}
\newcommand{\placeholderfigure}[2][0.95\linewidth]{%
  \begin{center}
    \fbox{\parbox[c][#2][c]{#1}{\centering Placeholder figure}}
  \end{center}
}

\begin{document}

\begin{frame}[t]
\begin{columns}[t]

\begin{column}{\sepwidth}\end{column}

% =============================
% Column 1
% =============================
\begin{column}{\colwidth}

\begin{alertblock}{Key Message}
\justifying
State the poster's main takeaway in 2--4 sentences. This should be the first thing a reader understands from several feet away.
\end{alertblock}

\begin{block}{Introduction}
\justifying
Briefly explain the problem, why it matters, and what gap your work addresses.

\begin{itemize}
  \item Motivation point one
  \item Motivation point two
  \item Why existing methods are limited
\end{itemize}
\end{block}

\begin{block}{Problem Setup}
\justifying
Define the setting, assumptions, or notation.

\begin{equation}
\mathcal{L}(\theta) = \sum_{i=1}^{N} \ell(f_{\theta}(x_i), y_i)
\end{equation}

\begin{itemize}
  \item Input:
  \item Output:
  \item Goal:
\end{itemize}
\end{block}

\begin{block}{Method Overview}
\justifying
Summarize your method at a high level. Keep this section scannable.

\placeholderfigure{10cm}

\begin{enumerate}
  \item Step one
  \item Step two
  \item Step three
\end{enumerate}
\end{block}

\end{column}

\begin{column}{\sepwidth}\end{column}

% =============================
% Column 2
% =============================
\begin{column}{\colwidth}

\begin{block}{Model / Approach}
\justifying
Describe the core idea more precisely.

\begin{equation}
p(z \mid x) \propto p(x \mid z)\,p(z)
\end{equation}

\begin{equation}
\mathrm{Score}(z, x) = -\log p(z) - \log p(x \mid z)
\end{equation}

\begin{itemize}
  \item What is optimized
  \item What is fixed
  \item Why the objective is useful
\end{itemize}
\end{block}

\begin{block}{Experimental Setup}
\justifying
List datasets, baselines, metrics, and implementation details.

\begin{itemize}
  \item Datasets:
  \item Baselines:
  \item Metrics:
  \item Training details:
\end{itemize}
\end{block}

\begin{alertblock}{Main Result}
\justifying
Highlight the most important quantitative or qualitative finding here.

\placeholderfigure{12cm}
\end{alertblock}

\begin{block}{Results Table}
\begin{center}
\begin{tabular}{lccc}
\toprule
Method & Metric 1 $\downarrow$ & Metric 2 $\uparrow$ & Time \\
\midrule
Baseline A & 0.81 & 72.4 & 1.0x \\
Baseline B & 0.74 & 75.9 & 1.3x \\
Ours       & \textbf{0.62} & \textbf{81.7} & 1.1x \\
\bottomrule
\end{tabular}
\end{center}
\end{block}

\end{column}

\begin{column}{\sepwidth}\end{column}

% =============================
% Column 3
% =============================
\begin{column}{\colwidth}

\begin{block}{Analysis}
\justifying
Use this section for ablations, interpretation, or error analysis.

\begin{itemize}
  \item Ablation one
  \item Ablation two
  \item Failure mode or caveat
\end{itemize}

\placeholderfigure{9cm}
\end{block}

\begin{block}{Discussion}
\justifying
Explain what the results mean and why they matter.

\begin{itemize}
  \item Practical implication
  \item Scientific implication
  \item Limitation
\end{itemize}
\end{block}

\begin{block}{Conclusion}
\justifying
Summarize the contribution in a few lines.

\begin{itemize}
  \item Main contribution
  \item Main empirical finding
  \item Future direction
\end{itemize}
\end{block}

\begin{block}{References}
\small
[1] Author et al. Title. Venue, Year.\\
[2] Author et al. Title. Venue, Year.\\
[3] Author et al. Title. Venue, Year.
\end{block}

\begin{alertblock}{Contact}
\small
Name Surname\\
email@university.edu\\
lab / website / QR code
\end{alertblock}

\end{column}

\begin{column}{\sepwidth}\end{column}

\end{columns}
\end{frame}

\end{document}
